\documentclass[11pt,a4paper]{article}

\usepackage[margin=1.5cm, columnsep=.8cm]{geometry}
\usepackage[onehalfspacing]{setspace}
\usepackage{graphicx}
\usepackage[breaklinks=true,colorlinks=true,linkcolor=blue,urlcolor=blue,citecolor=blue]{hyperref}
\usepackage{amsmath,amsthm,amssymb}
\usepackage{icomma}
\usepackage[indonesian]{babel}
\usepackage[T1]{fontenc}
\usepackage[numbers,sort&compress]{natbib}

\input{config.tex}

\linespread{1}

\begin{document}
\title{\includegraphics[width=5cm]{images/logo.png} \\ Spotify Wrapped: Seberapa \\ Pasaran Selera Musik Kamu?}
\author{%
Kelompok \group\\[.8em]
\begin{tabular}{@{}l r@{}}
\authorname{1} & (\authorid{1}) \\
\authorname{2} & (\authorid{2}) \\
\authorname{3} & (\authorid{3}) \\
\authorname{4} & (\authorid{4}) \\
\end{tabular}
}
\date{\today}
\maketitle
\vfill
\renewcommand\abstractname{Abstrak}
\section*{\abstractname}
\textit{If you only had 2 minutes to explain the work to a colleague, what should he/she know?}
\thispagestyle{empty}

\tableofcontents
\thispagestyle{empty}
\cleardoublepage
\pagenumbering{arabic} 
\newpage

\twocolumn

\section{Pendahuluan}
Musik merupakan bagian penting dari kehidupan sehari-hari, tiap
individu pastinya memiliki selera musik yang berbeda, contohnya
tiap orang mempunyai ketertarikan terhadap genre musik tertentu
\cite{han2022universality, dunn2012toward} seperti \textit{pop},
\textit{rock}, \textit{rnb}, dan banyak lagi. Tidak hanya dibatasi
dengan sekedar genre tertentu, kita sebagai makhluk sosial
cenderung suka mengikuti tren yang sedang populer, hal ini juga
melekat pada bagaimana cara kita mendengarkan musik. Oleh karena
itu, diperlukan visualisasi untuk memahami seberapa ``pasaran''
(\textit{mediocre}) selera musik seseorang dalam rentang waktu
tertentu.

Berbagai metode visualisasi telah dikembangkan untuk menampilkan
pola preferensi musik \cite{han2022universality}, contohnya mulai
dari grafik (\textit{chart}) \textit{Bubble}, \textit{Polar},
\textit{Bar}, dan banyak lagi. Akan tetapi, visualisasi ini biasanya
bersifat statis sehingga memberi kesan monoton \cite{bako2023user}.
Selain itu, hanya sedikit kostumisasi yang bisa dilakukan untuk
menyesuaikan grafik sesuai kebutuhan; misalnya ketika ingin membuat
antarmuka web, format output sering terbatas pada gambar (png, jpeg).
Hal ini dapat menimbulkan kerepotan karena visualisasi harus digambar
berulang kali.

Maka dari itu, kami merakit aplikasi visualisasi berupa poster interaktif
dengan antarmuka web menggunakan pustaka \texttt{d3} \cite{bostock2011d3}
sebagai alat bantu untuk menggambar grafik secara otomatis. Dengan memanfaatkan
teknologi web, kita dapat merakit alat tambahan (seperti telemetri, opsi
real-time). Hal ini dapat memudahkan pengembang dalam menyesuaikan visualisasi
berdasarkan umpan balik oleh pengguna dengan tepat dan efisien.

\section{Metodologi}
\textit{What are the important physical principles?}\\

\section{Eksperimen}
\textit{What would someone need to reproduce the experiment?}\\

\subsection{Dataset}
Dataset yang digunakan pada studi ini diambil secara dinamis dari API Spotify.

\subsection{Metrik Evaluasi}
Untuk mengukur efektivitas desain yang kami gunakan, maka hasil
telemetri $\Phi$ digunakan sebagai tolak ukur pembanding. Sumbu $x$
pada desain tidak secara langsung mendeskripsikan bahwa $x$ merupakan
popularitas (\%), tetapi dalam desain interaktif kami, akan muncul
deskripsi daripada sumbu $x$, maka untuk menguji efektifnya
metode ini, kita akan mencatat apakah pengguna $p$ berhasil (unfold) label
sumbu $x$, \ldots kita hitung Rata-rata Kursor Terarah (RKT) relatif
terhadap sumbu $x$:
\begin{equation}
% RKT = \frac{1}{|P|} \sum_{p \in P}{|\{i \mid i \in p \cap \Phi\}|}
% RKT = \frac{1}{|P|} \sum_{p \in P} \Phi_p
RKT = \frac{1}{|P|} \sum_{p \in P}{p \cap \Phi}
% RKT = \frac{1}{|P|} \sum_{p \in P} \Phi_p
% RKT = \frac{1}{|P|} \sum_{p \in P} |p \cap \Phi|
% RKT = \frac{1}{|P|} \sum_{p \in P} \mathbf{1}_{\{p \in \Phi\}}
\end{equation}

\section{Hasil}
\textit{What is (are) the measured answer(s) to the main question(s)?}\\

\section{Kesimpulan}
\textit{What is the final answer and how reliable is it? How could the experiment have been done differently or better?}\\

\renewcommand{\refname}{Referensi}
\bibliographystyle{IEEEtran}
\bibliography{report} 

\end{document}
